\section{Deterministic Placement Details}
\label{app:placement-details}

\subsection{Hierarchy Construction}

For each historical transaction, form a hyperedge \(E_t=\{s:a_t(s)=1\}\); edges with fewer than two scopes do not guide coarsening. Let \(\pi_\ell(E_t)\) be the distinct search groups touched by that edge at level \(\ell\). One fixed affinity rule for distinct groups \(G,H\) is
\begin{equation}
 A_\ell(G,H)=
 \sum_{\substack{t\in\mathcal H\\G,H\in\pi_\ell(E_t)}}
 \frac{1}{|\pi_\ell(E_t)|-1}.
\label{eq:coarsening-affinity}
\end{equation}
At each level, visit groups in canonical membership order and match an unmatched group with its highest-positive-affinity unmatched neighbor, subject to the bundle-size bound; break ties by canonical membership. Contract disjoint matched pairs in the search representation, carry unmatched groups forward, and record the hierarchy. Stop at the configured coarse-size target, level limit, or absence of an eligible match. Exact rational comparisons use bounded checked arithmetic; exceeding a construction bound makes the update ineligible, not machine-dependent.

\subsection{Canonical Output and Work Bounds}

To encode the incumbent, traverse hierarchy roots in canonical order. Emit each maximal node whose original scopes have one incumbent destination and satisfy the group-size bound; otherwise recurse into its children. This yields disjoint output groups without changing original-scope positions or routing weights.

The configuration bounds history size, declarations, retained scopes, hyperedge incidences, expanded metadata, hierarchy depth, refinement steps, candidate scans, and output bytes. Work accounting includes every affected transaction route, producer-load term, and dependency pair. Exhaustion cannot authorize a partially evaluated candidate; final improvement and feasibility checks still apply.

\subsection{Placement Readiness Subject}

The validation-and-custody subject binds the message domain, chain and validator epoch, update opportunity, old and new map digests and versions, finalized-history anchor, algorithm/schema configuration, and allowed activation rule. Signatures on different subjects cannot be combined. Retention and garbage collection preserve recovery for every still-applicable map and outstanding old-map block.

\subsection{Connectivity Does Not Determine Routing Cost}

For \(T_1=\{a,b\}\) and \(T_2=\{a,c\}\), with unit access weights and lowest-position tie-breaking, assignments \((a,b,c)=(1,2,3)\) and \((3,1,2)\) give the same hypergraph connectivity cost. Yet their transaction routes are respectively \((1,1)\) and \((1,2)\). If the operations on \(a\) are sensitive, their separated-dependency costs differ. A standard hypergraph connectivity score therefore cannot replace the routing-aware objective in the main paper.
